\chapter{复合玻色子、自旋表示与狄拉克倍增}

\section{交换对称性与复合玻色子}

开篇的问题是关于复合粒子的一个真正的悖论。如果两个费米子组成一个玻色子，那么两个这样的玻色子能否占据同一状态而不违反费米子排斥？答案要从交换规则本身说起。

假设两个全同粒子由含两个自变量的波函数描述。符号\(x\)和\(y\)标记两个粒子槽位；它们并不区分两个物种。对于玻色子，波函数是对称的，而对于费米子，波函数是反对称的：
\begin{equation}
\Psi_B(x,y)=\Psi_B(y,x),
\qquad
\Psi_F(x,y)=-\Psi_F(y,x).
\end{equation}
这种对称性表明，交换这两个不可区分粒子不会产生新的物理状态。

两个玻色子当然可以占据同一个单粒子态。如果该态的波函数为\(\varphi\)，则两粒子乘积为
\begin{equation}
\Psi(x,y)=\varphi(x)\varphi(y).
\end{equation}
波函数是普通的复值函数，而不是算符或矩阵，因此它们的值可交换：
\begin{equation}
\varphi(x)\varphi(y)=\varphi(y)\varphi(x).
\end{equation}
因此该乘积是对称的。同样的构造不可能给出非零的两费米子态，因为反对称性会要求
\begin{equation}
\varphi(x)\varphi(y)=-\varphi(y)\varphi(x),
\end{equation}
而不带负号的两边是相等的。

更一般地，从任意一个两粒子波函数\(F(x,y)\)出发，不假定任何交换对称性。其对称部分和反对称部分（忽略归一化）为
\begin{align}
F_{\mathrm{sym}}(x,y)&=F(x,y)+F(y,x),\\
F_{\mathrm{asym}}(x,y)&=F(x,y)-F(y,x).
\end{align}
交换检验是显式的：
\begin{align}
F_{\mathrm{sym}}(y,x)
&=F(y,x)+F(x,y)
=F_{\mathrm{sym}}(x,y),\\
F_{\mathrm{asym}}(y,x)
&=F(y,x)-F(x,y)
=-F_{\mathrm{asym}}(x,y).
\end{align}
加上交换后的函数会使其对称化；减去交换后的函数会使其反对称化。

一对氢原子提供了更有趣的情形。有四个组成费米子：两个全同电子和两个全同质子。令\(x_1,x_2\)表示电子变量，\(y_1,y_2\)表示质子变量。四体波函数必须在交换电子时以及在交换质子时分别反对称：
\begin{align}
\Psi(x_2,x_1;y_1,y_2)
&=-\Psi(x_1,x_2;y_1,y_2),\\
\Psi(x_1,x_2;y_2,y_1)
&=-\Psi(x_1,x_2;y_1,y_2).
\end{align}

为了引出两个不同的单原子波函数，设想一个氢原子定域在原点附近，波函数为\(\psi(x,y)\)，另一个定域在别处，波函数为\(\phi(x,y)\)。第二个波函数可以只是第一个的平移副本，即通过平移其电子和质子自变量得到，但在此设定中它仍然是一个不同的波函数。一个朴素的乘积是
\begin{equation}
\psi(x_1,y_1)\phi(x_2,y_2).
\end{equation}
它还不具有所要求的任何一种费米子反对称性。

首先对电子变量做反对称化：
\begin{equation}
\Psi_x=
\psi(x_1,y_1)\phi(x_2,y_2)
-\psi(x_2,y_1)\phi(x_1,y_2).
\end{equation}
在\(x_1\leftrightarrow x_2\)下，两项互换位置，整个表达式变号。现在再对质子变量对这个结果做反对称化。四项的排序由两次相继的反对称化固定。\footnote{这一构造还通过口头计算过程中所做的变量标签修正，固定了符号模式\(+,-,-,+\)。}
\begin{align}
\Psi_{2H}(x_1,x_2;y_1,y_2)
&=\Psi_x(x_1,x_2;y_1,y_2)\nonumber\\
&\quad{}-\Psi_x(x_1,x_2;y_2,y_1)\\
&=\psi(x_1,y_1)\phi(x_2,y_2)\nonumber\\
&\quad{}-\psi(x_2,y_1)\phi(x_1,y_2)\nonumber\\
&\quad{}-\psi(x_1,y_2)\phi(x_2,y_1)\nonumber\\
&\quad{}+\psi(x_2,y_2)\phi(x_1,y_1).
\label{eq:two-hydrogen-antisymmetrized}
\end{align}
交换电子会互换前两项和后两项，并伴随整体符号反转。交换质子也是如此。因此该态恰好具有组成费米子所要求的两种独立反对称性。

\begin{classroomqa}
\lecturetimestamp{00:11:11}
\audiencequestion{如果氢原子是玻色子，那么是什么阻止两个氢原子被直接放在彼此之上？}
\lecturerresponse{玻色子交换统计并不会消除原子之间的力。组成粒子之间的泡利排斥，连同其他原子间力，可以使空间重叠的代价极其高昂。这与询问原子是否可以占据同一离域动量态是不同的。动量本征态弥散于整个空间，因此共同动量并不意味着把原子撞到同一点上。甚至像台球这样的硬核物体也可以服从玻色子交换统计：它们的硬核会阻止空间重叠，而玻色子统计仍然允许共同动量占据。}
\end{classroomqa}

\begin{classroomqa}
\lecturetimestamp{00:12:28}
\audiencequestion{把两个氢原子置于同一动量本征态，是否会把它们的组成费米子置于同一状态？}
\lecturerresponse{不会。四体态在交换两个电子变量时仍然反对称，并且在交换两个质子变量时也分别反对称。只有同时交换两种组成粒子---即交换两个完整原子---才是对称的。}
\end{classroomqa}

确实，交换两个原子意味着同时进行两次替换
\begin{equation}
x_1\leftrightarrow x_2,
\qquad
y_1\leftrightarrow y_2.
\end{equation}
第一次替换贡献一个费米子负号，第二次替换再贡献一个：
\begin{equation}
\begin{aligned}
\Psi_{2H}(x_2,x_1;y_2,y_1)
&=(-1)(-1)\Psi_{2H}(x_1,x_2;y_1,y_2)\\
&=\Psi_{2H}(x_1,x_2;y_1,y_2).
\end{aligned}
\end{equation}
同样的结果可以在式~\eqref{eq:two-hydrogen-antisymmetrized}中逐项看出。在整体原子交换下，
\begin{equation}
\psi(x_1,y_1)\phi(x_2,y_2)
\longleftrightarrow
\psi(x_2,y_2)\phi(x_1,y_1),
\end{equation}
因此第一项和第四项互换。同样，
\begin{equation}
-\psi(x_2,y_1)\phi(x_1,y_2)
\longleftrightarrow
-\psi(x_1,y_2)\phi(x_2,y_1),
\end{equation}
因此第二项和第三项互换。完整状态在原子互换下是对称的。两个组成负号使复合体系获得正号。

这一交换陈述在逻辑上独立于使原子重叠所需的能量。强排斥可能使这种重叠无法实现，但这并不改变对称交换定律。

\begin{classroomqa}
\lecturetimestamp{00:14:42}
\audiencequestion{所构造的四体态在物理上代表什么，它是否具体是动量空间中的态？}
\lecturerresponse{它是由两个单原子波函数先对电子标签反对称化、再对质子标签反对称化而得到的双原子波函数。它不必写在动量空间中。这些变量可以描述位置、动量或其他表示；所需的组成粒子交换对称性相同。}
\end{classroomqa}

因此，两个原子可以占据同一个离域动量波函数，尽管迫使它们的空间分布强烈重叠可能付出很大的能量代价。

现在令两个单原子波函数相等，\(\phi=\psi\)。式~\eqref{eq:two-hydrogen-antisymmetrized}变为
\begin{align}
\Psi_{\mathrm{same}}
&=\psi(x_1,y_1)\psi(x_2,y_2)
-\psi(x_2,y_1)\psi(x_1,y_2)\nonumber\\
&\quad
-\psi(x_1,y_2)\psi(x_2,y_1)
+\psi(x_2,y_2)\psi(x_1,y_1).
\end{align}

这个四体态会消失吗？不会。第一项和第四项乘积相等，第二项和第三项乘积相等。利用单原子波函数的交换性，
\begin{equation}
\Psi_{\mathrm{same}}
=
2\left[
\psi(x_1,y_1)\psi(x_2,y_2)
-\psi(x_2,y_1)\psi(x_1,y_2)
\right],
\label{eq:same-hydrogen-state}
\end{equation}
它一般不为零。它在两个电子变量下仍然反对称，并且分别在两个质子变量下反对称。

这里并没有把某个定域明确的电子分配给某个定域明确的原子。束缚在原子中的电子在其原子波函数上空间离域。因此，让两个原子具有相同的原子态并不会使\(x_1=x_2\)成立，也不会使\(y_1=y_2\)成立。

组成粒子的排除效应可以通过直接检验重合来看到。如果两个模型化电子自变量重合，\(x_1=x_2=x\)，式~\eqref{eq:same-hydrogen-state}给出
\begin{align}
\Psi_{\mathrm{same}}(x,x;y_1,y_2)
&=2\bigl[
\psi(x,y_1)\psi(x,y_2)\nonumber\\
&\quad{}-\psi(x,y_1)\psi(x,y_2)
\bigr]\\
&=0.
\end{align}
当\(y_1=y_2\)时，相应的表达式也为零。因此，同态复合波函数不为零，而它在相同费米子自变量重合处的振幅为零。\footnote{在这个模型中，自旋以及任何其他单粒子标签都被省略。费米子反对称性适用于完整的单粒子标签集合；这里展示的抵消涉及此处保留的自变量重合。}

\begin{classroomqa}
\lecturetimestamp{00:20:41}
\audiencequestion{原子能否处于同一状态，而它们的电子和质子并不处于相同的组成粒子状态？}
\lecturerresponse{可以。复合原子波函数相等并不会消除电子和质子各自的反对称性。完整的四体态可以描述两个原子处于同一个原子态，同时对相同费米子自变量重合赋予零振幅。}
\end{classroomqa}

\begin{classroomqa}
\lecturetimestamp{00:21:12}
\audiencequestion{反对称化是否投影掉了相同组成粒子重合的那些部分？}
\lecturerresponse{是的。在此处展示的变量中，那些重合分量精确抵消。精确的陈述是：当两个相同费米子自变量相等时，反对称波函数为零。}
\end{classroomqa}

这种抵消并不依赖于先把原子波函数设为相等。对于式~\eqref{eq:two-hydrogen-antisymmetrized}中的一般态，
\begin{align}
\Psi_{2H}(x,x;y_1,y_2)
&=\psi(x,y_1)\phi(x,y_2)\nonumber\\
&\quad{}-\psi(x,y_1)\phi(x,y_2)\nonumber\\
&\quad{}-\psi(x,y_2)\phi(x,y_1)\nonumber\\
&\quad{}+\psi(x,y_2)\phi(x,y_1)\\
&=0.
\end{align}

\begin{classroomqa}
\lecturetimestamp{00:21:32}
\audiencequestion{当原子处于不同状态且相同组成粒子变量重合时，这种抵消仍然成立吗？}
\lecturerresponse{是的。它来自反对称性本身，而不是来自令\(\phi=\psi\)。同样的陈述在另一种表示中也成立：变换到动量空间后，当两个相同费米子动量自变量重合时，只要其他被省略的单粒子标签也重合，振幅就为零。}
\end{classroomqa}

\section{作为有限维量子系统的自旋}

自旋的数学在同位旋、色以及粒子物理的其他内部对称性中反复出现。这种类比是形式上的：出现的是同一套代数，尽管这些物理量不必具有相同的解释。当前的任务是把自旋再稍微展开一些，然后回到狄拉克方程。这套工具之后将用于同位旋，这是一个历史上与1930年代相关的概念。

电子是自旋为\(\tfrac12\)的粒子，质子和{\proseunicodefont μ}子也是如此。其自旋的\(z\)分量有两个可能取值。在设定单位之前，写出
\begin{equation}
S_z=m\hbar.
\end{equation}
从现在起\(\hbar=1\)。字母\(L\)按惯例保留给轨道角动量，因此内禀自旋将用\(S\)表示。它遵循相同的角动量代数：
\begin{equation}
[S_x,S_y]=iS_z,
\quad
[S_y,S_z]=iS_x,
\quad
[S_z,S_x]=iS_y,
\end{equation}
以及
\begin{equation}
S_\pm=S_x\pm iS_y.
\end{equation}
算符\(S_+\)升高\(S_z\)本征值，而\(S_-\)降低它。

\begin{classroomqa}
\lecturetimestamp{00:25:39}
\audiencequestion{讨论是否专门针对内禀自旋？}
\lecturerresponse{内禀自旋是预期的焦点。相同的角动量代数更普遍地适用，包括轨道角动量或宏观物体的转动。}
\end{classroomqa}

对于自旋二分之一，唯一的\(S_z\)本征值是\(+\tfrac12\)和\(-\tfrac12\)。暂时抑制关于电子的其他一切：忘记它的位置和动量，只保留两个自旋态。于是
\begin{equation}
|\psi\rangle
=
\alpha|\uparrow_z\rangle+\beta|\downarrow_z\rangle
\quad\longleftrightarrow\quad
\begin{pmatrix}
\alpha\\
\beta
\end{pmatrix},
\end{equation}
其中
\begin{equation}
|\alpha|^2=P(\uparrow_z),
\qquad
|\beta|^2=P(\downarrow_z),
\qquad
|\alpha|^2+|\beta|^2=1.
\end{equation}
这是一个二维态空间。这并不是说普通空间是二维的。

\begin{classroomqa}
\lecturetimestamp{00:28:18}
\audiencequestion{是否说过物理空间是二维的？}
\lecturerresponse{没有。自旋态空间是二维的。位置和动量已被暂时抑制，就好像电子被钉在墙上，只能测量其自旋。}
\end{classroomqa}

这个态空间上的算符由\(2\times2\)矩阵表示。方便的是选择\(S_z\)基，使\(S_z\)为对角，但对角性是一种基的选择，而不是物理要求。

\begin{classroomqa}
\lecturetimestamp{00:29:46}
\audiencequestion{\(S_z\)必须是对角的吗？}
\lecturerresponse{不。选择它对角是为了方便。其他矩阵表示在物理上等价。}
\end{classroomqa}

\begin{classroomqa}
\lecturetimestamp{00:30:08}
\audiencequestion{\(S_z\)是一个矩阵吗？}
\lecturerresponse{\(S_z\)是一个可观测量，因此是一个算符。在所选有限维基中，该算符由作用于态矢量的矩阵表示。}
\end{classroomqa}

一个特别有用的表示是
\begin{equation}
S_i=\frac12\sigma_i,
\end{equation}
其中
\begin{equation}
\sigma_x=
\begin{pmatrix}
0&1\\
1&0
\end{pmatrix},
\quad
\sigma_y=
\begin{pmatrix}
0&-i\\
i&0
\end{pmatrix},
\quad
\sigma_z=
\begin{pmatrix}
1&0\\
0&-1
\end{pmatrix}.
\end{equation}
不带二分之一因子的矩阵是泡利矩阵；自旋算符本身包含该因子。直接相乘可验证对易关系。该表示并不唯一：旋转坐标轴或改变基会产生另一组等价的三个矩阵。

这是自旋二分之一角动量的二分量表示。

\subsection{沿坐标轴的本征态}

逐分量求解本征值方程是有用的。由于\(\sigma_i\)的本征值是\(\pm1\)，\(S_i=\sigma_i/2\)的本征值是\(\pm\tfrac12\)。

对于\(\sigma_z\)，两个方程是
\begin{align}
\begin{pmatrix}
1&0\\
0&-1
\end{pmatrix}
\begin{pmatrix}
\alpha\\
\beta
\end{pmatrix}
&=
\begin{pmatrix}
\alpha\\
\beta
\end{pmatrix},
&
\begin{pmatrix}
1&0\\
0&-1
\end{pmatrix}
\begin{pmatrix}
\alpha\\
\beta
\end{pmatrix}
&=
-
\begin{pmatrix}
\alpha\\
\beta
\end{pmatrix}.
\end{align}

\begin{figure}[t]
\centering
\includegraphics[width=0.74\textwidth]{lecture_08_figure_02.png}
\caption{\(+1\)本征值方程，其中\(\sigma_z\)作用于二分量旋量。}
\lectureframe{00:34:57}
\end{figure}

对于本征值\(+1\)，逐行相乘给出
\begin{equation}
\begin{pmatrix}
1\cdot\alpha+0\cdot\beta\\
0\cdot\alpha-\beta
\end{pmatrix}
=
\begin{pmatrix}
\alpha\\
-\beta
\end{pmatrix}
=
\begin{pmatrix}
\alpha\\
\beta
\end{pmatrix}.
\end{equation}
上面的方程自动满足，而下面的方程给出\(-\beta=\beta\)，因此\(\beta=0\)。归一化要求\(|\alpha|=1\)；选择无关紧要的整体相位使\(\alpha=1\)，
\begin{equation}
|+z\rangle=
\begin{pmatrix}
1\\
0
\end{pmatrix}.
\end{equation}

对于本征值\(-1\)，
\begin{equation}
\begin{pmatrix}
\alpha\\
-\beta
\end{pmatrix}
=
\begin{pmatrix}
-\alpha\\
-\beta
\end{pmatrix}.
\end{equation}
现在下面的方程自动满足，上面的方程给出\(\alpha=-\alpha\)，所以\(\alpha=0\)。选择\(\beta=1\)给出
\begin{equation}
|-z\rangle=
\begin{pmatrix}
0\\
1
\end{pmatrix}.
\end{equation}

现在改变实验并测量\(S_x\)。由于
\begin{equation}
\sigma_x
\begin{pmatrix}
\alpha\\
\beta
\end{pmatrix}
=
\begin{pmatrix}
0\cdot\alpha+\beta\\
\alpha+0\cdot\beta
\end{pmatrix}
=
\begin{pmatrix}
\beta\\
\alpha
\end{pmatrix},
\end{equation}
\(+1\)方程变为
\begin{equation}
\begin{pmatrix}
\beta\\
\alpha
\end{pmatrix}
=
\begin{pmatrix}
\alpha\\
\beta
\end{pmatrix}.
\end{equation}
两个分量方程\(\beta=\alpha\)和\(\alpha=\beta\)是同一个条件。归一化给出
\begin{equation}
2|\alpha|^2=1,
\end{equation}
因此一个方便的选择是
\begin{equation}
|+x\rangle
=
\frac{1}{\sqrt2}
\begin{pmatrix}
1\\
1
\end{pmatrix}.
\end{equation}

\begin{figure}[t]
\centering
\includegraphics[width=0.74\textwidth]{lecture_08_figure_03.png}
\caption{对于本征值\(+1\)，\(\sigma_x\)交换两个分量，要求\(\alpha=\beta\)。}
\lectureframe{00:39:42}
\end{figure}

\begin{classroomqa}
\lecturetimestamp{00:40:13}
\audiencequestion{如果包含二分之一因子，本征值会是二分之一吗？}
\lecturerresponse{是的。泡利矩阵\(\sigma_x\)的本征值是\(\pm1\)，而\(S_x=\sigma_x/2\)的本征值是\(\pm\tfrac12\)。}
\end{classroomqa}

对于本征值\(-1\)，
\begin{equation}
\begin{pmatrix}
\beta\\
\alpha
\end{pmatrix}
=
-
\begin{pmatrix}
\alpha\\
\beta
\end{pmatrix},
\end{equation}
所以
\begin{equation}
\beta=-\alpha,
\qquad
\alpha=-\beta.
\end{equation}
这些同样是同一个方程。归一化后，
\begin{equation}
|-x\rangle
=
\frac{1}{\sqrt2}
\begin{pmatrix}
1\\
-1
\end{pmatrix}.
\end{equation}
因此，\(z\)向上和\(z\)向下的等量叠加沿正\(x\)轴极化，而带有相对负号的相同叠加沿负\(x\)轴极化。

对于\(y\)方向，
\begin{equation}
\sigma_y
\begin{pmatrix}
\alpha\\
\beta
\end{pmatrix}
=
\begin{pmatrix}
-i\beta\\
i\alpha
\end{pmatrix}.
\end{equation}
两个归一化本征矢量为
\begin{equation}
|+y\rangle
=
\frac{1}{\sqrt2}
\begin{pmatrix}
1\\
i
\end{pmatrix},
\qquad
|-y\rangle
=
\frac{1}{\sqrt2}
\begin{pmatrix}
1\\
-i
\end{pmatrix},
\end{equation}
在整体相位意义下成立。这两个态给出\(z\)向上概率\(1/2\)、\(z\)向下概率\(1/2\)。相对相位\(i\)在这两个概率中不可见，却区分了\(y\)极化与\(x\)极化。线性组合也描述沿任意方向极化的自旋；该计算留待后文。

\begin{classroomqa}
\lecturetimestamp{00:44:22}
\audiencequestion{\(i\)分量的符号是否被赋给了错误的\(\sigma_y\)本征值？}
\lecturerresponse{该符号确实可能被赋给了错误的本征值。此处这一赋值尚不确定；不凭记忆将其固定。}
\end{classroomqa}

直接检验为
\begin{align}
\sigma_y\frac{1}{\sqrt2}
\begin{pmatrix}
1\\
i
\end{pmatrix}
&=
\frac{1}{\sqrt2}
\begin{pmatrix}
-i\,i\\
i
\end{pmatrix}
=
\frac{1}{\sqrt2}
\begin{pmatrix}
1\\
i
\end{pmatrix},\\
\sigma_y\frac{1}{\sqrt2}
\begin{pmatrix}
1\\
-i
\end{pmatrix}
&=
\frac{1}{\sqrt2}
\begin{pmatrix}
(-i)(-i)\\
i
\end{pmatrix}
=
-\frac{1}{\sqrt2}
\begin{pmatrix}
1\\
-i
\end{pmatrix}.
\end{align}
因此\(+i\)属于本征值\(+1\)，而\(-i\)属于本征值\(-1\)。\footnote{口头给出的赋值在多次反转后仍未确定；此处显示的符号由所展示的矩阵乘法固定。} 每个自旋二分之一态，包括沿任意方向极化的每个态，都是\(|+z\rangle\)与\(|-z\rangle\)的线性组合。

\section{自旋一与位置相关旋量}

对于自旋\(1\)，可能的\(S_z\)取值为\(+1,0,-1\)。因此态空间是三维的。自旋\(1\)是整数自旋，对应玻色统计。我们寻找三个满足相同角动量对易关系的\(3\times3\)矩阵。

下面给出一个在三个坐标方向之间对称的表示。\footnote{\(S_y\)的符号遵循\(00{:}48{:}50\)附近可见的矩阵，并且也由\([S_x,S_y]=iS_z\)固定。反转该符号将反转所需的对易关系。}
\begin{equation}
\begin{aligned}
S_x
&=
i
\begin{pmatrix}
0&0&0\\
0&0&1\\
0&-1&0
\end{pmatrix},
&
S_y
&=
i
\begin{pmatrix}
0&0&1\\
0&0&0\\
-1&0&0
\end{pmatrix},\\
S_z
&=
i
\begin{pmatrix}
0&1&0\\
-1&0&0\\
0&0&0
\end{pmatrix}.
\end{aligned}
\end{equation}
每个矩阵都是\(i\)乘以一个实反对称矩阵。每个矩阵都包含一个\(1\)、一个\(-1\)，以及与自身坐标相关的零行和零列：第一个对应\(S_x\)，第二个对应\(S_y\)，第三个对应\(S_z\)。这给出一个有用的助记法。在此表示中\(S_z\)不是对角的。

令
\begin{equation}
v=
\begin{pmatrix}
\alpha\\
\beta\\
\gamma
\end{pmatrix}.
\end{equation}
则
\begin{equation}
S_zv
=
\begin{pmatrix}
i\beta\\
-i\alpha\\
0
\end{pmatrix}.
\end{equation}
对于本征值\(+1\)，
\begin{equation}
i\beta=\alpha,
\qquad
-i\alpha=\beta,
\qquad
0=\gamma.
\end{equation}
因此\(\beta=-i\alpha\)且\(\gamma=0\)。归一化给出\(2|\alpha|^2=1\)，所以
\begin{equation}
|m=+1\rangle
=
\frac{1}{\sqrt2}
\begin{pmatrix}
1\\
-i\\
0
\end{pmatrix}.
\end{equation}
对于本征值\(-1\)，
\begin{equation}
i\beta=-\alpha,
\qquad
-i\alpha=-\beta,
\qquad
0=-\gamma,
\end{equation}
因此\(\beta=i\alpha\)、\(\gamma=0\)，并且
\begin{equation}
|m=-1\rangle
=
\frac{1}{\sqrt2}
\begin{pmatrix}
1\\
i\\
0
\end{pmatrix}.
\end{equation}
最后，零本征值方程给出\(\alpha=\beta=0\)，留下\(\gamma\)任意。其归一化解为
\begin{equation}
|m=0\rangle
=
\begin{pmatrix}
0\\
0\\
1
\end{pmatrix}.
\end{equation}
没有其他本征值：这三个本征矢量穷尽了三维空间。因子\(1/\sqrt2\)并非装饰性的。当该代数应用于同位旋时，同样的系数后来会进入定量实验预言。

\begin{classroomqa}
\lecturetimestamp{00:54:10}
\audiencequestion{这些自旋一本征矢量的复分量应如何解读？}
\lecturerresponse{对于\(m=+1\)，分量为\(1/\sqrt2,-i/\sqrt2,0\)；对于\(m=-1\)，分量为\(1/\sqrt2,i/\sqrt2,0\)；对于\(m=0\)，分量为\(0,0,1\)。}
\end{classroomqa}

其他矩阵三元组满足同一代数，并且在物理上等价。上述对称表示之所以有用，是因为三个方向被同等对待，而且它还会再次被使用。

\subsection{恢复位置}

现在恢复粒子的运动。自旋与位置和动量都对易，因此位置或动量中的任一个都可以与一个自旋分量同时指定。位置与动量不对易，而不同的自旋分量彼此也不对易。

一个运动的自旋二分之一粒子由位置相关旋量描述，
\begin{equation}
\psi(x)=
\begin{pmatrix}
\psi_\uparrow(x)\\
\psi_\downarrow(x)
\end{pmatrix}.
\end{equation}
量
\begin{equation}
\rho_{\uparrow_z}(x)=|\psi_\uparrow(x)|^2,
\qquad
\rho_{\downarrow_z}(x)=|\psi_\downarrow(x)|^2
\end{equation}
是联合的位置与自旋概率密度。它们的和
\begin{equation}
\rho(x)=\psi^\dagger(x)\psi(x)
=|\psi_\uparrow(x)|^2+|\psi_\downarrow(x)|^2,
\end{equation}
是位置密度。如果已知粒子在\(x\)处被探测到，则条件概率为
\begin{equation}
P(\uparrow_z\mid x)
=
\frac{|\psi_\uparrow(x)|^2}{\psi^\dagger(x)\psi(x)},
\qquad
P(\downarrow_z\mid x)
=
\frac{|\psi_\downarrow(x)|^2}{\psi^\dagger(x)\psi(x)},
\end{equation}
只要\(\rho(x)\neq0\)。因此电子由二分量的空间波函数表示，其中局域自旋振幅可以随位置变化。

对于运动的自旋一粒子，写
\begin{equation}
\phi(x)=
\begin{pmatrix}
\phi_1(x)\\
\phi_2(x)\\
\phi_3(x)
\end{pmatrix}.
\end{equation}
在空间每一点，三个分量构成上述同一自旋空间中的一个矢量。自旋一粒子具有整数自旋，是玻色子。

\begin{classroomqa}
\lecturetimestamp{01:00:10}
\audiencequestion{\(\phi_1(x)\)本身是自旋一结果\(m=+1\)的概率吗？}
\lecturerresponse{不是。它是在所选三分量基中的振幅。要得到确定的自旋结果，需构造相应归一化本征矢量的左矢，将场投影到其上，并取模平方。}
\end{classroomqa}

\begin{figure}[t]
\centering
\includegraphics[width=0.82\textwidth]{lecture_08_figure_04.png}
\caption{归一化的\(\sigma_y\)本征矢量与三个自旋二分之一矩阵并列。}
\lectureframe{01:00:12}
\end{figure}

例如，
\begin{equation}
\begin{aligned}
|m=+1\rangle
&=
\frac{1}{\sqrt2}
\begin{pmatrix}
1\\
-i\\
0
\end{pmatrix}
\\
\Longrightarrow\quad\langle m=+1|
&=
\frac{1}{\sqrt2}
\begin{pmatrix}
1&i&0
\end{pmatrix}.
\end{aligned}
\end{equation}
因此投影振幅为
\begin{align}
\langle m=+1|\phi(x)\rangle
&=
\frac{1}{\sqrt2}
\begin{pmatrix}
1&i&0
\end{pmatrix}
\begin{pmatrix}
\phi_1(x)\\
\phi_2(x)\\
\phi_3(x)
\end{pmatrix}\\
&=
\frac{\phi_1(x)+i\phi_2(x)}{\sqrt2}.
\end{align}
将同一规则应用于全部三个本征矢量，便得到自旋分辨的密度。\footnote{这些公式在明确修正了\(\phi_1\)与\(\phi_2\)不必为实数之后，完成了投影代数。若把它们当作实数，就会错误地抹去区分\(m=+1\)与\(m=-1\)密度的干涉项。}
\begin{align}
\rho_{+1}(x)
&=
\frac12\left|\phi_1(x)+i\phi_2(x)\right|^2,\\
\rho_{-1}(x)
&=
\frac12\left|\phi_1(x)-i\phi_2(x)\right|^2,\\
\rho_0(x)
&=
|\phi_3(x)|^2.
\end{align}
它们满足
\begin{equation}
\begin{aligned}
\rho_{+1}(x)+\rho_{-1}(x)+\rho_0(x)
&=
|\phi_1(x)|^2+|\phi_2(x)|^2
\\
&\quad{}+|\phi_3(x)|^2
\\
&=
\phi^\dagger(x)\phi(x).
\end{aligned}
\end{equation}
若粒子被条件化为已在\(x\)处被探测到，则
\begin{equation}
P(m\mid x)
=
\frac{\rho_m(x)}{\phi^\dagger(x)\phi(x)},
\end{equation}
只要\(\phi^\dagger(x)\phi(x)\neq0\)。一般测量规则始终相同：取所需归一化本征矢量的左矢，将实际态投影到其上，再取模平方。所得密度保留其位置依赖性。

更高自旋的矩阵更为复杂，但原理不变。

\begin{classroomqa}
\lecturetimestamp{01:07:31}
\audiencequestion{自旋三分之二的状态空间有多大？}
\lecturerresponse{它有四个状态，其中\(m=\tfrac32,\tfrac12,-\tfrac12,-\tfrac32\)。自旋二有五个状态。}
\end{classroomqa}

\begin{classroomqa}
\lecturetimestamp{01:07:55}
\audiencequestion{为什么讨论涉及了\(x,y,z\)却没有时间方向？这一处理是非相对论性的吗？}
\lecturerresponse{到目前为止它是非相对论性的。自旋与质量一样，最简便的描述是在粒子的静止系中给出的；它在另一惯性系中的描述由相应的洛伦兹变换得到。这一相对论性步骤留待以后处理。}
\end{classroomqa}

狄拉克方程提供了相对论性框架，并解释了为什么其四分量波函数包含一个自旋双重态。

\section{狄拉克矩阵与四分量旋量}

四分量狄拉克波函数满足
\begin{equation}
i\frac{\partial\Psi}{\partial t}
=
-i\alpha_i\frac{\partial\Psi}{\partial x_i}
+m\beta\Psi,
\qquad i=1,2,3,
\end{equation}
其中重复的空间指标被求和。因此导数项意味着
\begin{equation}
\alpha_1\frac{\partial\Psi}{\partial x}
+\alpha_2\frac{\partial\Psi}{\partial y}
+\alpha_3\frac{\partial\Psi}{\partial z}.
\end{equation}
对于平面波，这变为代数方程
\begin{equation}
\omega\Psi=(\alpha_i k_i+m\beta)\Psi.
\end{equation}

\begin{classroomqa}
\lecturetimestamp{01:11:32}
\audiencequestion{\(m\)的两次出现表示同一事物吗？}
\lecturerresponse{不。一次是被用作空间指标，另一次是粒子的质量。将空间指标重命名为\(i\)；当\(i\)出现两次时，它从\(1\)到\(3\)求和。}
\end{classroomqa}

对频率算符取平方必须重现相对论关系
\begin{equation}
\omega^2=\mathbf{k}^2+m^2.
\end{equation}
因此平方项要求，对每个\(i\)，
\begin{equation}
\alpha_i^2=I_4,
\qquad
\beta^2=I_4.
\end{equation}

\begin{classroomqa}
\lecturetimestamp{01:13:01}
\audiencequestion{\(\beta\)对\(\Psi\)的四个分量中的每一个都不同吗？}
\lecturerresponse{不。只有一个矩阵\(\beta\)，作用于整个四分量旋量，并且它满足\(\beta^2=I_4\)。}
\end{classroomqa}

混合项必须消失。对于不同的空间方向，以及每个\(\alpha_i\)与\(\beta\)之间，这给出
\begin{equation}
\alpha_i\alpha_j+\alpha_j\alpha_i=0
\quad (i\neq j),
\qquad
\alpha_i\beta+\beta\alpha_i=0.
\end{equation}
等价地，
\begin{equation}
\{\alpha_i,\alpha_j\}=2\delta_{ij}I_4,
\quad
\{\alpha_i,\beta\}=0,
\quad
\{\beta,\beta\}=2I_4.
\end{equation}
每个矩阵与每一个不同的矩阵反交换，而它与自身的反交换子是单位矩阵的两倍。

如果没有\(\beta\)，三个泡利矩阵就已经提供了三个平方为单位矩阵且两两反交换的\(2\times2\)矩阵。质量项要求第四个这样的矩阵。不存在具有所需性质的第四个\(2\times2\)矩阵，而\(3\times3\)矩阵也不能解出该代数。因此最小的表示是\(4\times4\)。

一个\(4\times4\)矩阵可以分成四个\(2\times2\)块。在一种方便的表示中，
\begin{equation}
\alpha_i=
\begin{pmatrix}
\sigma_i&0\\
0&-\sigma_i
\end{pmatrix},
\qquad i=1,2,3.
\end{equation}

\begin{figure}[t]
\centering
\includegraphics[width=0.68\textwidth]{lecture_08_figure_05.png}
\caption{用\(2\times2\)块写出的狄拉克矩阵\(\alpha_i\)，其中\(\sigma_i\)与\(-\sigma_i\)位于对角线上。}
\lectureframe{01:17:23}
\end{figure}

\begin{classroomqa}
\lecturetimestamp{01:17:41}
\audiencequestion{在较早的表示中，\(\sigma_i\)块是以另一种方式排列的吗？}
\lecturerresponse{是的。在另一种表示中它们被放在非对角线上。这两种排列是等价的，并由相似变换相联系。}
\end{classroomqa}

在目前的表示中，剩余的矩阵是
\begin{equation}
\beta=
\begin{pmatrix}
0&I_2\\
I_2&0
\end{pmatrix},
\qquad
I_2=
\begin{pmatrix}
1&0\\
0&1
\end{pmatrix}.
\end{equation}
块乘法把对\(4\times4\)代数的检验归结为\(2\times2\)矩阵的普通乘法。

\begin{classroomqa}
\lecturetimestamp{01:19:17}
\audiencequestion{较早的表示对\(\beta\)使用了不同的符号或对角形式吗？}
\lecturerresponse{它使用了不同但物理上等价的基。基变换混合\(\Psi\)的四个分量，并改变矩阵的外观而不改变物理。}
\end{classroomqa}

\(\Psi\) 的四个分量并不是时空的四个坐标。它们是在这个反对易代数能够作用的最小向量空间中的坐标。这里有两种不同的二重结构：分为 \(2\times2\) 块的划分，以及每个块内部的两分量泡利结构。预先考虑下文所确立的解释，这四种可能性可以组织为
\begin{equation}
\begin{array}{c|cc}
& S_z=+\tfrac12 & S_z=-\tfrac12\\ \hline
\omega>0 & (\omega>0,+\tfrac12) & (\omega>0,-\tfrac12)\\
\omega<0 & (\omega<0,+\tfrac12) & (\omega<0,-\tfrac12)
\end{array}.
\end{equation}
下面一行由负频电子态组成。它不是正电子行；正电子将被解释为这些态被填满后的一个空穴。

\section{静止系能量加倍与正电子}

二重能量结构在零动量处最容易显露。当 \(\mathbf{k}=0\) 时，平面波态在空间上是常数，空间导数项消失。静止系方程\footnote{在 \(01{:}23{:}38\) 处的参考系在质量项上保留了一个瞬时的额外因子 \(i\)；在 \(01{:}24{:}05\) 处的口头更正将其去掉。} 为
\begin{equation}
i\frac{\partial\Psi}{\partial t}=m\beta\Psi.
\end{equation}

\begin{figure}[t]
\centering
\includegraphics[width=0.76\textwidth]{lecture_08_figure_06.png}
\caption{在零动量处，空间导数消失，留下修正后的静止系方程 \(i\partial_t\Psi=m\beta\Psi\)。}
\lectureframe{01:23:38}
\end{figure}

将四分量列向量分为上、下两个两分量对象：
\begin{equation}
\Psi=
\begin{pmatrix}
\psi_+\\
\psi_-
\end{pmatrix}.
\end{equation}

\begin{classroomqa}
\lecturetimestamp{01:24:47}
\audiencequestion{标签 \(\psi_+\) 和 \(\psi_-\) 已经意味着正能量和负能量了吗？}
\lecturerresponse{不。在这个阶段，它们只是上、下两个两分量块的名称。它们的能量组合尚未被求出。}
\end{classroomqa}

\begin{classroomqa}
\lecturetimestamp{01:25:10}
\audiencequestion{被划分的四分量列向量代表什么？}
\lecturerresponse{它就是同一个四分量狄拉克旋量，只是被组织成与矩阵的块形式相匹配。每一半都是一个两分量对象，而两个这样的对象构成原来的四分量列向量。}
\end{classroomqa}

矩阵 \(\beta\) 交换这两个块：
\begin{equation}
\beta
\begin{pmatrix}
\psi_+\\
\psi_-
\end{pmatrix}
=
\begin{pmatrix}
\psi_-\\
\psi_+
\end{pmatrix}.
\end{equation}
因此静止系方程变为
\begin{equation}
i\frac{\partial\psi_+}{\partial t}=m\psi_-,
\qquad
i\frac{\partial\psi_-}{\partial t}=m\psi_+.
\end{equation}

有两种自然的建议。再取一次时间导数会使系统解耦，但将两个方程相加和相减会立即将它们对角化。确实，
\begin{align}
i\frac{\partial}{\partial t}(\psi_++\psi_-)
&=m(\psi_++\psi_-),\\
i\frac{\partial}{\partial t}(\psi_+-\psi_-)
&=-m(\psi_+-\psi_-).
\end{align}
定义
\begin{equation}
\chi_+=\psi_++\psi_-,
\qquad
\chi_-=\psi_+-\psi_-,
\end{equation}
得到两个独立的方程，
\begin{equation}
i\partial_t\chi_+=m\chi_+,
\qquad
i\partial_t\chi_-=-m\chi_-.
\end{equation}
对于确定的频率依赖关系 \(e^{-i\omega t}\)，和具有
\begin{equation}
\omega=+m,
\end{equation}
而差具有
\begin{equation}
\omega=-m.
\end{equation}
这正是先前一维例子中发现的正、负能量加倍：相对论方程再次具有正能量和负能量的电子解。因此，块二重结构分离为正、负能量分支。纯静止系态可以写为
\begin{equation}
\Psi_{\omega=+m}=
\begin{pmatrix}
\xi\\
\xi
\end{pmatrix},
\qquad
\Psi_{\omega=-m}=
\begin{pmatrix}
\xi\\
-\xi
\end{pmatrix},
\end{equation}
其中 \(\xi\) 仍然是一个任意的两分量对象。

和与差已经处理掉了一种二重结构，但 \(\xi\) 的两个分量仍然存在。泡利矩阵作用在这个剩余的二重对上。这就是与自旋相关的二重对，它受下文讨论的旋转和守恒定律检验的约束。

\begin{classroomqa}
\lecturetimestamp{01:32:23}
\audiencequestion{自旋可以被当作外层的二重结构，而能量符号被当作内层的二重结构吗？}
\lecturerresponse{可以。两种因子排序是等价的。人们可以交换哪个指标写在前面；这一改变只是四个分量的重新排列。}
\end{classroomqa}

负频分支导致狄拉克海解释。负频率被赋予负的电子能量。填满每一个负能量电子态会在泡利不相容原理允许的范围内尽可能降低能量，因为每个完整的费米子态只能容纳一个电子。完全填满的负能量海于是被称为真空。

将这个填满的海中的一个电子激发到正能态。结果包含一个正能电子和一个未被占据的负能态。移除一个负能量会提高能量，因此空穴具有正的激发能量；由于缺少一个电子，它还带有相反的电荷。这个空穴就是正电子。

\begin{classroomqa}
\lecturetimestamp{01:32:57}
\audiencequestion{正电子解释可以再说明一次吗？}
\lecturerresponse{两个块的和具有正频率，而差具有负频率。负频解是负能量电子态，而真空通过填满所有这些态得到。将一个已填满的电子激发到正能量会同时留下一个正能电子和一个空穴。那个具有正能量、带相反电荷的空穴就是正电子。}
\end{classroomqa}

因此，代数同时产生了两个后果：正能支和负能支，它们成为粒子--反粒子结构；以及一个与自旋二分之一相关的残余二分量结构。引人注目的是，这两者都来自能够满足相对论性反对易关系的最小矩阵。当这一数学构造产生出带自旋的粒子以及一个反粒子支时，自旋作为经验性质才刚刚出现不久；在那个历史背景下，这一结果是戏剧性的。

\begin{classroomqa}
\lecturetimestamp{01:35:14}
\audiencequestion{负能电子本身就是正电子吗？}
\lecturerresponse{不是。负能电子态在真空中是被填满的。正电子是从那个被填满的集合中缺少一个电子：是一个空穴，而不是一个负能电子。}
\end{classroomqa}

空穴可以参与跃迁。正电子具有正能，因为它是缺少一个负能电子。如果一个正能电子填入空穴，电子和正电子就会湮灭，但它们的能量必须被带走，通常由光子带走。用海的语言说，一个电子填满了一个未被占据的负能态。同样的能量释放可以描绘为原子中一个激发态电子落到较低能级：向下跃迁发射光子。

\begin{classroomqa}
\lecturetimestamp{01:36:39}
\audiencequestion{为什么更高的负能电子不会落入空穴？}
\lecturerresponse{这个问题混淆了两种情况。一个正能电子填入空穴是电子--正电子湮灭，光子带走能量。负能海中的一个电子发生跃迁则是另一种过程，下面会讨论。}
\end{classroomqa}

仅零动量计算本身，在把能量符号分离之后，已经确立了一个二重简并。两个态的简并本身并不证明这些态带有角动量。

\begin{classroomqa}
\lecturetimestamp{01:37:44}
\audiencequestion{零动量计算如何证明剩下的那一对态是自旋？}
\lecturerresponse{这一对态可以称为上和下，泡利矩阵也出现在完整方程中，但仅凭这一点还不是证明。必须证明这一对态在空间旋转下如何变换，或者证明赋予它自旋能使总角动量在相互作用和碰撞中守恒。这一论证留待以后进行。}
\end{classroomqa}

对于正能静止态，上块和下块相等，
\begin{equation}
\Psi_{\omega=+m}=
\begin{pmatrix}
\xi\\
\xi
\end{pmatrix}.
\end{equation}
确立\(\xi\)是一个角动量二重态，需要刚才描述的旋转论证或守恒论证；它并不能仅仅由数出两个分量而得出。

\begin{classroomqa}
\lecturetimestamp{01:40:46}
\audiencequestion{另一个电子即使能量不同，也能与这个正电子湮灭吗？}
\lecturerresponse{能。电子不保留个体身份。任一可用的电子都可以填入空穴，而初始能量和动量的差异则反映在发射光子的数目和能量上。}
\end{classroomqa}

空穴图像也解释了完全发生在负能海内部的跃迁。假设一个被占据的负能态移入空穴。原来的空穴被填满，但那个电子原来所在的态现在空了。空穴移动了。等价地说，正电子改变了它的能量并发射了辐射。

\begin{classroomqa}
\lecturetimestamp{01:42:12}
\audiencequestion{为什么邻近的负能能级不能相继落入空穴？}
\lecturerresponse{它们可以把空穴从一个能级移到另一个能级。一个电子落入旧空穴，会在它原先的能量处留下一个新空穴，因此初始正电子变成能量不同的正电子。能量差由光子带走。}
\end{classroomqa}

\begin{classroomqa}
\lecturetimestamp{01:44:39}
\audiencequestion{为什么这不会产生无休止的向下辐射级联？}
\lecturerresponse{因为能量和动量必须同时守恒。这两条守恒定律不允许这种不受限制的过程。普通自由电子也会出现同样的问题：它不能仅仅通过在真空中发射光子来不断降低自己的能量。}
\end{classroomqa}
